\documentclass[11pt]{article}
\usepackage[margin=1in]{geometry}
\usepackage{amsmath,amssymb,amsthm,mathtools}
\usepackage[T1]{fontenc}
\usepackage[utf8]{inputenc}
\usepackage{enumitem}
\usepackage{booktabs,longtable,array}
\usepackage{hyperref}
\usepackage{xcolor}
\usepackage{microtype}
\hypersetup{colorlinks=true, linkcolor=blue!50!black, citecolor=blue!50!black, urlcolor=blue!50!black}

\newtheorem{theorem}{Theorem}[section]
\newtheorem{proposition}[theorem]{Proposition}
\newtheorem{lemma}[theorem]{Lemma}
\newtheorem{corollary}[theorem]{Corollary}
\newtheorem{claim}[theorem]{Claim}
\theoremstyle{definition}
\newtheorem{definition}[theorem]{Definition}
\newtheorem{assumption}[theorem]{Assumption}
\theoremstyle{remark}
\newtheorem{remark}[theorem]{Remark}

\newcommand{\KL}{\mathrm{KL}}
\newcommand{\Ent}{\mathrm{Ent}}
\newcommand{\Prob}{\mathsf{P}}
\newcommand{\E}{\mathbb{E}}
\newcommand{\R}{\mathbb{R}}
\newcommand{\N}{\mathbb{N}}
\newcommand{\Pcal}{\mathcal{P}}
\newcommand{\Mcal}{\mathcal{M}}
\newcommand{\Qcal}{\mathcal{Q}}
\newcommand{\Zcal}{\mathcal{Z}}
\newcommand{\Fcal}{\mathcal{F}}
\newcommand{\Gcal}{\mathcal{G}}
\newcommand{\W}{\mathcal{W}}
\newcommand{\one}{\mathbf{1}}
\newcommand{\dd}{\,d}
\newcommand{\supp}{\operatorname{supp}}
\newcommand{\argmin}{\operatorname*{arg\,min}}
\newcommand{\argmax}{\operatorname*{arg\,max}}
\newcommand{\esssup}{\operatorname*{ess\,sup}}
\newcommand{\ri}{\operatorname{ri}}
\newcommand{\dom}{\operatorname{dom}}
\newcommand{\Law}{\operatorname{Law}}
\newcommand{\Var}{\operatorname{Var}}
\newcommand{\Cov}{\operatorname{Cov}}
\newcommand{\diag}{\operatorname{diag}}
\newcommand{\Id}{\mathrm{Id}}
\newcommand{\Lip}{\operatorname{Lip}}
\newcommand{\osc}{\operatorname{osc}}
\newcommand{\diam}{\operatorname{diam}}
\newcommand{\proj}{\operatorname{proj}}
\newcommand{\e}{\mathrm{e}}

\newenvironment{argument}{\paragraph{Argument.}\begin{enumerate}[leftmargin=2em]}{\end{enumerate}}
\newenvironment{proofoutline}{\paragraph{Proof architecture.}\begin{enumerate}[leftmargin=2em]}{\end{enumerate}}


% Source-faithfulness apparatus used by the paper reconstructions.
\newcommand{\sourceanchor}[1]{\par\smallskip\noindent\textbf{Source anchor.} #1\par\smallskip}
\newcommand{\sourcecomment}[1]{\begin{remark}#1\end{remark}}
\newenvironment{sourceguide}
  {\begin{description}[style=nextline,leftmargin=3.2cm,labelwidth=3cm]}
  {\end{description}}

% Backward-compatible environments used by some of the older reconstructions.
\newenvironment{distilled}[1][]{\begin{theorem}[#1]}{\end{theorem}}
\newenvironment{distillation}{\begin{enumerate}[leftmargin=2em]}{\end{enumerate}}

\newenvironment{commentarynotes}{\begin{remark}\begin{enumerate}[leftmargin=2em]}{\end{enumerate}\end{remark}}

\title{Stochastic Control Liaisons: Schr\"odinger Bridges, Entropic Transport, and Sinkhorn Iteration}
\author{}
\date{Chen, Georgiou, and Pavon (2021)}
\begin{document}
\maketitle

\section*{Source map and scope}
\begin{sloppypar}
\begin{sourceguide}
\item[Primary source] \path{CGP2021-2005.10963v3.pdf}, checked against its arXiv source archive.
\item[Target chain] The dynamic/static Schr\"odinger equivalence, factorization into Schr\"odinger potentials, Girsanov/control-energy formulation, generalized HJB/Hopf--Cole verification theorem, regularized optimal transport, continuous Fortet formulation, and the finite Fortet--IPF--Sinkhorn existence/uniqueness/convergence theorem.
\item[Source anchors] Sections 5.2--5.5, equations labeled \texttt{decomposition}, \texttt{staticindex}, \texttt{representation}, \texttt{optimaljoint}, \texttt{opt1}--\texttt{BConestep}, \texttt{girsanov}, \texttt{eq:regularizedOT}; Section 6, \texttt{FDProblem}, \texttt{Lagrangian2}, \texttt{prioroptcond}, \texttt{HJB}, \texttt{generalizedschrodinger}, and the following theorem; Sections 7.1 and 7.3; Sections 8.1--8.4, Theorem \texttt{Fundtheorem}, Lemmas \texttt{Linear}, \texttt{inversion}, and \texttt{compothm}.
\end{sourceguide}
\end{sloppypar}

This note reconstructs the central mathematical chain in the order used by the paper.  Historical
material, examples, and applications are omitted unless they enter a proof.  Where the source
states a result by citation or gives only a formal variational calculation, that status is retained
rather than replaced silently by a different proof.

\section{Path-space formulation and entropy decomposition}
Let $\Omega=C([0,1];\R^n)$ with canonical process $X_t$.  Let $W$ denote a reference Markov path
measure, initially Wiener measure, and let $\mathcal D(\rho_0,\rho_1)$ be the set of path measures
whose time-zero and time-one marginals have densities $\rho_0$ and $\rho_1$.
The Schr\"odinger bridge problem is
\begin{equation}
 P_{\rm SBP}
 =\argmin\{\KL(P\|W):P\in\mathcal D(\rho_0,\rho_1)\}.
 \label{eq:CGP-bridge}
\end{equation}
The paper motivates this by Sanov's theorem and the Gibbs conditioning principle, but the proof
chain begins from the entropy minimization itself.

For a finite-energy diffusion $P$, disintegrate both $P$ and $W$ with respect to
$(X_0,X_1)$:
\[
 P(d\omega)=P_{xy}(d\omega)\rho^P_{01}(dx,dy),
 \qquad
 W(d\omega)=W_{xy}(d\omega)\rho^W_{01}(dx,dy).
\]
The relative-entropy chain rule gives the source equation \texttt{decomposition}:
\begin{align}
 \KL(P\|W)
 &=\KL(\rho^P_{01}\|\rho^W_{01})
 +\int\KL(P_{xy}\|W_{xy})\,\rho^P_{01}(dx,dy). \label{eq:CGP-decomp}
\end{align}
The endpoint law and conditional bridge laws can be varied independently.  The second term is
nonnegative and vanishes exactly when $P_{xy}=W_{xy}$ for $\rho^P_{01}$-almost every $(x,y)$.
Consequently the dynamic problem reduces to
\begin{equation}
 \inf_{\pi\in\Pi(\rho_0,\rho_1)}\KL(\pi\|\rho^W_{01}).
 \label{eq:CGP-static}
\end{equation}
If $\pi^*$ solves this static problem, then
\begin{equation}
 P^*(\cdot)=\int W_{xy}(\cdot)\,\pi^*(dx,dy)
 \label{eq:CGP-representation}
\end{equation}
solves the path-space problem.  This law remains in the reciprocal class of the reference: its
three-point bridge transition is the same as that of $W$.

\begin{remark}
The source explicitly notes that \eqref{eq:CGP-representation} remains valid for a possibly
non-Markovian finite-energy prior.  The later multiplicative transition formula and drift formula,
however, use Markov structure.
\end{remark}

\section{Static entropy minimization and the Schr\"odinger system}
Assume the reference endpoint law has density
\[
 \rho^W_{01}(x,y)=\rho^W_0(x)p(0,x;1,y),
\]
with strictly positive transition density $p$.  Introduce Lagrange multipliers $\lambda(x)$ and
$\mu(y)$ for the two marginals.  The Lagrangian is
\begin{align*}
 \mathcal L(\rho_{01},\lambda,\mu)
 &=\iint \rho_{01}\log\frac{\rho_{01}}{\rho^W_{01}}\,dxdy\\
 &\quad+\int\lambda(x)\left(\int\rho_{01}(x,y)\,dy-\rho_0(x)\right)dx\\
 &\quad+\int\mu(y)\left(\int\rho_{01}(x,y)\,dx-\rho_1(y)\right)dy.
\end{align*}
The first variation in $\rho_{01}$ is
\[
 1+\log\rho^*_{01}(x,y)-\log p(0,x;1,y)-\log\rho^W_0(x)
 +\lambda(x)+\mu(y)=0.
\]
Exponentiating separates the $x$ and $y$ variables.  With positive functions
$\widehat\varphi(x)$ and $\varphi(y)$,
\begin{equation}
 \rho^*_{01}(x,y)
 =\widehat\varphi(x)p(0,x;1,y)\varphi(y).
 \label{eq:CGP-optimaljoint}
\end{equation}
The marginal equations are
\begin{align}
 \widehat\varphi(x)\int p(0,x;1,y)\varphi(y)\,dy&=\rho_0(x),
 \label{eq:CGP-opt1}\\
 \varphi(y)\int p(0,x;1,y)\widehat\varphi(x)\,dx&=\rho_1(y).
 \label{eq:CGP-opt2}
\end{align}
Define propagated potentials
\begin{align*}
 \widehat\varphi(1,y)&=\int p(0,x;1,y)\widehat\varphi(0,x)\,dx,\\
 \varphi(0,x)&=\int p(0,x;1,y)\varphi(1,y)\,dy.
\end{align*}
Then the system is the pair of linear propagation equations coupled by
\begin{equation}
 \varphi(0,x)\widehat\varphi(0,x)=\rho_0(x),
 \qquad
 \varphi(1,y)\widehat\varphi(1,y)=\rho_1(y).
 \label{eq:CGP-boundary-products}
\end{equation}
The potentials are determined only up to
$(\varphi,\widehat\varphi)\mapsto(a\varphi,a^{-1}\widehat\varphi)$.
The marginal flow factors as
\[
 \rho(t,x)=\varphi(t,x)\widehat\varphi(t,x).
\]
The source cites Fortet, Beurling, Jamison, and F\"ollmer for existence and uniqueness in the
continuous setting~\cite{Fortet1940,Beurling1960,Jamison1974,Jamison1975,Follmer1988}; it does not
prove that theorem here.

\section{Brownian reference and entropic optimal transport}
Let $W_\varepsilon$ be Brownian motion with variance parameter $\varepsilon$, so
\[
 p_\varepsilon(0,x;1,y)
 =(2\pi\varepsilon)^{-n/2}
 \exp\left(-\frac{\|x-y\|^2}{2\varepsilon}\right).
\]
For a coupling $\pi$ with fixed marginals,
\begin{align*}
 \KL(\pi\|\rho^{W_\varepsilon}_{01})
 &=\frac1{2\varepsilon}\iint\|x-y\|^2\,d\pi
   -S(\pi)+C(\rho_0,W_\varepsilon),
\end{align*}
where the last term is independent of $\pi$.  Multiplication by $\varepsilon$ shows that the static
bridge problem is equivalent to
\begin{equation}
 \min_{\pi\in\Pi(\rho_0,\rho_1)}
 \left\{\frac12\iint\|x-y\|^2\,d\pi
 +\varepsilon\iint \rho_{01}(x,y)\log\rho_{01}(x,y)\,dxdy\right\}.
 \label{eq:CGP-regOT}
\end{equation}
This is the paper's precise sense in which the Schr\"odinger problem is an entropy-regularized
optimal-transport problem.

\section{Girsanov and stochastic control}
For a Brownian reference with diffusion coefficient $\sqrt\varepsilon$, let a controlled law satisfy
\[
 dX_t=u_t\,dt+\sqrt\varepsilon\,dW_t,
 \qquad X_0\sim\rho_0,
 \qquad X_1\sim\rho_1.
\]
The source invokes Girsanov's formula~\cite{Girsanov1960,KaratzasShreve1988}
\begin{equation}
 \KL(P^u\|W)
 =\KL((P^u)_0\|W_0)
 +\E_{P^u}\int_0^1\frac{\|u_t\|^2}{2\varepsilon}\,dt.
 \label{eq:CGP-girsanov}
\end{equation}
With the initial marginal fixed, the first term is constant.  Hence the bridge is equivalent to the
minimum-energy steering problem
\begin{equation}
 \inf_u\E\int_0^1\frac{\|u_t\|^2}{2\varepsilon}\,dt
 \quad\text{subject to the controlled SDE and endpoint laws.}
 \label{eq:CGP-stoch-control}
\end{equation}
The optimal feedback is
\begin{equation}
 u^*(t,x)=\varepsilon\nabla\log\varphi(t,x),
 \label{eq:CGP-brownian-feedback}
\end{equation}
where
\begin{align*}
 \partial_t\varphi+\frac\varepsilon2\Delta\varphi&=0,\\
 \partial_t\widehat\varphi-\frac\varepsilon2\Delta\widehat\varphi&=0,
\end{align*}
with the endpoint product conditions \eqref{eq:CGP-boundary-products}.

The corresponding density/control formulation is
\begin{align}
 \inf_{\rho,u}&\int_0^1\int\frac12\|u\|^2\rho\,dxdt,
 \label{eq:CGP-fluid-cost}\\
 \partial_t\rho+\nabla\cdot(u\rho)-\frac\varepsilon2\Delta\rho&=0,
 \qquad \rho(0)=\rho_0,\ \rho(1)=\rho_1. \label{eq:CGP-fluid-fp}
\end{align}
The source cautions that this fluid problem reproduces the optimal marginal flow but not the full
path law: temporal correlations and path regularity are lost.

\section{General Markov prior}
For
\[
 dX_t=f(t,X_t)\,dt+\sqrt\varepsilon\,dW_t,
\]
the same argument yields
\[
 dX_t=[f(t,X_t)+u_t]dt+\sqrt\varepsilon\,dW_t,
 \qquad
 u_t^*=\varepsilon\nabla\log\varphi(t,X_t),
\]
with
\begin{align*}
 \partial_t\varphi+f\cdot\nabla\varphi+\frac\varepsilon2\Delta\varphi&=0,\\
 \partial_t\widehat\varphi+\nabla\cdot(f\widehat\varphi)
 -\frac\varepsilon2\Delta\widehat\varphi&=0.
\end{align*}
The transition densities satisfy the multiplicative-functional relation
\begin{equation}
 p^*(s,x;t,y)
 =\frac{1}{\varphi(s,x)}\bar p(s,x;t,y)\varphi(t,y).
 \label{eq:CGP-multiplicative}
\end{equation}
This relation, not merely the bridge mixture \eqref{eq:CGP-representation}, gives the modified
Markov dynamics.

\section{Generalized control problem with killing or creation}
\sourceanchor{Section 6, equations \texttt{Fokker-Planck}, \texttt{FDProblem}, \texttt{Lagrangian2}, \texttt{prioroptcond}, \texttt{HJB}, and the theorem following \texttt{generalizedschrodinger}.}
Let $a=\sigma\sigma^\top$ be positive semidefinite of constant rank.  The prior density equation is
\begin{equation}
 \partial_t\rho+\nabla\cdot(f\rho)+V\rho
 =\frac12\sum_{i,j}\partial_{ij}(a_{ij}\rho).
 \label{eq:CGP-prior-fp}
\end{equation}
The state potential $V\ge0$ allows loss of mass.  The controlled problem is
\begin{align}
 \inf_{\rho,u}&\int_0^1\int
 \left(\frac12\|u\|^2+V\right)\rho\,dxdt,
 \label{eq:CGP-FD-cost}\\
 \partial_t\rho+\nabla\cdot((f+\sigma u)\rho)
 &=\frac12\sum_{i,j}\partial_{ij}(a_{ij}\rho),
 \label{eq:CGP-FD-constraint}\\
 \rho(0)&=\rho_0,\qquad \rho(1)=\rho_1.
\end{align}

Introduce a smooth Lagrange multiplier $\lambda$.  After integration by parts and deletion of
boundary terms fixed by the endpoint constraints, the integrand becomes
\[
 \left[\frac12\|u\|^2+V-
 \left(\partial_t\lambda+(f+\sigma u)\cdot\nabla\lambda
 +\frac12\sum_{i,j}a_{ij}\partial_{ij}\lambda\right)\right]\rho.
\]
Pointwise completion of the square in $u$ gives
\begin{equation}
 u^*=\sigma^\top\nabla\lambda. \label{eq:CGP-u-star}
\end{equation}
The coefficient of $\rho$ then vanishes when
\begin{equation}
 \partial_t\lambda+f\cdot\nabla\lambda
 +\frac12\sum_{i,j}a_{ij}\partial_{ij}\lambda
 +\frac12\nabla\lambda\cdot a\nabla\lambda-V=0.
 \label{eq:CGP-HJB}
\end{equation}
The controlled density equation becomes
\begin{equation}
 \partial_t\rho+\nabla\cdot((f+a\nabla\lambda)\rho)
 =\frac12\sum_{i,j}\partial_{ij}(a_{ij}\rho).
 \label{eq:CGP-opt-evolution}
\end{equation}

\begin{proposition}[Source verification proposition]
If a smooth $\lambda$ solves \eqref{eq:CGP-HJB} and a density $\rho^*$ solves
\eqref{eq:CGP-opt-evolution} with the prescribed endpoints, then
$u^*=\sigma^\top\nabla\lambda$ and $\rho^*$ solve the generalized control problem.
\end{proposition}
\begin{proof}
For arbitrary $u$,
\[
 \frac12\|u\|^2-u\cdot\sigma^\top\nabla\lambda
 =\frac12\|u-\sigma^\top\nabla\lambda\|^2
 -\frac12\nabla\lambda\cdot a\nabla\lambda.
\]
Substitute this into the integrated Lagrangian and use \eqref{eq:CGP-HJB}.  All terms independent
of the square are fixed by the endpoint conditions, while the remaining square is nonnegative and
vanishes exactly at $u^*$.
\end{proof}

Set $\varphi=e^\lambda$.  Direct differentiation gives
\[
 \partial_t\varphi=\varphi\partial_t\lambda,
 \quad
 \partial_{ij}\varphi
 =\varphi(\partial_{ij}\lambda+
          \partial_i\lambda\partial_j\lambda).
\]
Hence \eqref{eq:CGP-HJB} becomes the linear backward equation
\begin{equation}
 \partial_t\varphi+f\cdot\nabla\varphi
 +\frac12\sum_{i,j}a_{ij}\partial_{ij}\varphi=V\varphi.
 \label{eq:CGP-harmonic}
\end{equation}
Define $\widehat\varphi=\rho/\varphi$.  Substitution into
\eqref{eq:CGP-opt-evolution} and cancellation with \eqref{eq:CGP-harmonic} gives
\begin{equation}
 \partial_t\widehat\varphi+\nabla\cdot(f\widehat\varphi)
 -\frac12\sum_{i,j}\partial_{ij}(a_{ij}\widehat\varphi)
 =-V\widehat\varphi.
 \label{eq:CGP-adjoint}
\end{equation}
Together with
\[
 \varphi(0)\widehat\varphi(0)=\rho_0,
 \qquad
 \varphi(1)\widehat\varphi(1)=\rho_1,
\]
these are the generalized Schr\"odinger equations.  Thus a positive solution gives
\[
 u^*=\sigma^\top\nabla\log\varphi,
 \qquad \rho^*=\varphi\widehat\varphi.
\]
The source states this as a theorem.  Existence and uniqueness of the PDE system are not proved in
this section; they are referred to Beurling--Jamison theory under positivity of the fundamental
solution.

\section{Free energy and the discrete entropic transport problem}
For a finite state space with transport costs $c_{ij}$ and marginals $p,q$, the regularized problem is
\begin{equation}
 \inf_{\pi\in\Pi(p,q)}
 \left[\sum_{i,j}c_{ij}\pi_{ij}-\varepsilon S(\pi)\right],
 \qquad
 S(\pi)=-\sum_{i,j}\pi_{ij}\log\pi_{ij}.
 \label{eq:CGP-discrete-regOT}
\end{equation}
If $g_{ij}=e^{-c_{ij}/\varepsilon}$, this is equivalent, up to marginal-dependent constants, to
minimizing relative entropy from the Gibbs matrix $G=(g_{ij})$.  The row/column scaling equations are therefore the finite Schr\"odinger system; the corresponding
alternating normalization procedure is the Sinkhorn iteration~\cite{Sinkhorn1964,SinkhornKnopp1967}.

\section{Continuous Fortet equation}
\sourceanchor{Section 8.1, equations \texttt{OPT1}, \texttt{OPT2}, and \texttt{nonlineareq}.}
The paper states Fortet's theorem rather than reproving it.  Let $g\ge0$ be continuous and bounded,
positive outside null sections, and let continuous nonnegative $\rho_0,\rho_1$ have equal mass.
Assume the crucial integrability condition
\[
 \int\frac{\rho_1(y)}{\int g(z,y)\rho_0(z)\,dz}\,dy<\infty.
\]
Then
\begin{align*}
 \widehat\varphi(x)\int g(x,y)\varphi(y)\,dy&=\rho_0(x),\\
 \varphi(y)\int g(x,y)\widehat\varphi(x)\,dx&=\rho_1(y)
\end{align*}
has a solution with $\varphi$ continuous nonnegative and $\widehat\varphi$ measurable nonnegative,
with zeros only where the corresponding marginal vanishes.

Eliminating $\widehat\varphi$ gives the nonlinear fixed-point map
\begin{equation}
 \mathcal C(h)(x)
 =\int g(x,y)
 \frac{\rho_1(y)}{\int g(z,y)\rho_0(z)/h(z)\,dz}\,dy.
 \label{eq:CGP-Fortet-map}
\end{equation}
Fortet works on measurable functions bounded below by a positive constant and finite almost
everywhere.  The map is order-preserving and homogeneous of degree one, but it need not preserve
the positive lower bound.  The paper identifies this failure as the fundamental obstruction to a
straight Hilbert-metric contraction in the continuous function-space cone and refers to Fortet's three-sequence truncation scheme for the remedy~\cite{Fortet1940}.  Later
proof-level reconstructions of that scheme include Essid--Pavon and L\'eonard~\cite{EssidPavon2019,Leonard2019}.

\section{Finite Schr\"odinger system: statement and iteration}
\sourceanchor{Section 8.2, Theorem \texttt{Fundtheorem}, equations \texttt{eq:DSchroesystem}, \texttt{Composition}, and \texttt{Iteration}.}
Let $\mathcal X=\{1,\dots,n\}$, let $p,q$ be probability vectors, and let
$G=(g_{ij})$ have strictly positive entries.

\begin{theorem}[Fundamental finite theorem]\label{thm:CGP-finite}
There exist vectors
$\varphi_0,\widehat\varphi_0,\varphi_1,\widehat\varphi_1$ satisfying
\begin{align}
 \varphi_{0,i}&=\sum_jg_{ij}\varphi_{1,j},
 &\widehat\varphi_{1,j}&=\sum_ig_{ij}\widehat\varphi_{0,i},
 \label{eq:CGP-finite-linear}\\
 \varphi_{0,i}\widehat\varphi_{0,i}&=p_i,
 &\varphi_{1,j}\widehat\varphi_{1,j}&=q_j.
 \label{eq:CGP-finite-products}
\end{align}
The positive potential rays are unique up to reciprocal rescaling.
The Fortet--IPF--Sinkhorn iteration converges projectively to this solution ray.
\end{theorem}

Define
\[
 \mathcal E x=Gx,
 \qquad \mathcal E^\dagger x=G^\top x,
 \qquad \mathcal D_0x=p/x,
 \qquad \mathcal D_1x=q/x,
\]
with division componentwise, and
\begin{equation}
 \mathcal C=\mathcal E\circ\mathcal D_1\circ
              \mathcal E^\dagger\circ\mathcal D_0.
 \label{eq:CGP-composition}
\end{equation}
The source initializes $\varphi_0^{(0)}=\mathbf1$ and iterates
\begin{equation}
 \varphi_0^{(k+1)}=\mathcal C(\varphi_0^{(k)}).
 \label{eq:CGP-iteration}
\end{equation}
Strict positivity of $G$ ensures that the linear maps send every nonzero nonnegative vector to the
interior.  Thus the componentwise divisions appearing after those linear steps are defined even
when some entries of $p$ or $q$ vanish.

\section{Hilbert metric lemmas used in the finite proof}
For $x,y>0$, the Hilbert distance is
\[
 d_H(x,y)=\log\frac{\max_i x_i/y_i}{\min_i x_i/y_i}.
\]
It is a metric on rays.  The source recalls three external Birkhoff--Bushell results~\cite{Birkhoff1957,Bushell1973Projective,Bushell1973Hilbert}: monotone
homogeneous maps of degree $r$ are $r$-Lipschitz; a positive linear map has contraction coefficient
$\tanh(\Delta/4)$; and a strict contraction on a complete normalized cone section has a unique
fixed ray.  The use of this geometry for matrix scaling follows the Franklin--Lorenz line of argument~\cite{FranklinLorenz1989}.

\begin{lemma}[Source Lemma \texttt{Linear}]
The maps $\mathcal E$ and $\mathcal E^\dagger$ are strict Hilbert contractions.
\end{lemma}
\begin{proof}
For a positive matrix,
\[
 \Delta(\mathcal E)
 =\sup_{i,j,k,\ell}
 \log\frac{g_{ij}g_{k\ell}}{g_{i\ell}g_{kj}}<\infty.
\]
Birkhoff's formula gives
\[
 \kappa(\mathcal E)=\tanh(\Delta(\mathcal E)/4)<1,
\]
and the same argument applies to $G^\top$.
\end{proof}

\begin{lemma}[Source Lemma \texttt{inversion}]
$\mathcal D_0$ and $\mathcal D_1$ are Hilbert nonexpansive; when the marginal vector is strictly
positive, each is an isometry on the interior.
\end{lemma}
\begin{proof}
Componentwise inversion interchanges maxima and reciprocal minima:
\[
 d_H(x^{-1},y^{-1})=d_H(x,y).
\]
Componentwise multiplication by a positive fixed vector cancels in every ratio:
\[
 d_H(p\odot x,p\odot y)=d_H(x,y).
\]
If $p$ has zero coordinates, the source replaces equality by the corresponding restriction/inequality.
\end{proof}

\begin{lemma}[Source Lemma \texttt{compothm}]
The composition $\mathcal C$ is a strict contraction on positive rays.
\end{lemma}
\begin{proof}
Contraction coefficients multiply under composition.  The two nonlinear diagonal-inversion maps
have coefficient at most one, while both positive linear maps have coefficient strictly below one.
Therefore
\[
 \kappa(\mathcal C)
 \le\kappa(\mathcal E)\kappa(\mathcal E^\dagger)<1.
\]
\end{proof}

\section{Proof of the finite theorem}
The ray space of the interior positive orthant, equivalently any normalized positive simplex
section, is complete in Hilbert distance.  By the preceding lemma and Banach--Caccioppoli,
$\mathcal C$ has a unique fixed ray and the iteration \eqref{eq:CGP-iteration} converges to it
projectively.

Let $\varphi_0$ represent this ray.  Then for some $\alpha>0$,
\begin{equation}
 \mathcal C(\varphi_0)=\alpha\varphi_0.
 \label{eq:CGP-alpha-ray}
\end{equation}
Define successively
\[
 \widehat\varphi_0=p/\varphi_0,
 \qquad
 \widehat\varphi_1=G^\top\widehat\varphi_0,
 \qquad
 \varphi_1=q/\widehat\varphi_1.
\]
Then $G\varphi_1=\alpha\varphi_0$, and the product identities hold except that the initial one is
scaled by $\alpha$ when the source rewrites the cycle.  Pair with the Euclidean inner product:
\begin{align*}
 1
 &=\alpha\langle\widehat\varphi_0,\varphi_0\rangle\\
 &=\alpha\langle\widehat\varphi_0,G\varphi_1\rangle\\
 &=\alpha\langle G^\top\widehat\varphi_0,\varphi_1\rangle\\
 &=\alpha\langle\widehat\varphi_1,\varphi_1\rangle\\
 &=\alpha.
\end{align*}
The first and last inner products are the total masses of $p$ and $q$, both equal to one.  Thus
$\alpha=1$, so the fixed ray contains an actual fixed vector satisfying the four Schr\"odinger
equations.  Uniqueness of the ray follows from strict contraction.

The resulting coupling
\[
 \pi_{ij}=\widehat\varphi_{0,i}g_{ij}\varphi_{1,j}
\]
has row sums $p$ and column sums $q$ by \eqref{eq:CGP-finite-linear}--\eqref{eq:CGP-finite-products}.
It is the unique entropy-minimizing coupling because the finite relative-entropy objective is
strictly convex on the affine set of couplings.

\section{Source-faithfulness and technical comments}
\begin{commentarynotes}
\item The paper's dynamic/static reduction relies on the relative-entropy disintegration identity.
It is not proved from first principles in the survey; it is cited to F\"ollmer.
\item The first-variation derivation of the Schr\"odinger system is formal unless positivity,
absolute continuity, and differentiability of the entropy functional are justified.  The source uses
it as a sufficient optimality calculation and cites separate existence theory.
\item The generalized HJB theorem is a verification theorem: it assumes a sufficiently regular
positive solution of the coupled PDE system.  It does not prove PDE existence from the displayed
coefficients and endpoints.
\item In the source, the derivation of the equation for $\widehat\varphi=\rho/\varphi$ is described as
``a long but straightforward calculation.''  The reconstruction records the substitution and
cancellation but does not replace it by a probabilistic argument.
\item The continuous Fortet subsection is an outline of Fortet's theorem, not a reproduction of
Fortet's three-sequence proof.  That proof belongs to the separate Fortet and L\'eonard
reconstructions.
\item The finite proof establishes a unique fixed ray first and then proves the ray multiplier is one
from conservation of total mass.  This scalar argument is essential and must not be replaced by
assuming a normalized fixed point at the outset.
\item The source theorem permits probability vectors with zero coordinates while speaking of
positive potential vectors.  The surrounding text explains that strict positivity of $G$ keeps the
linear images positive and that the diagonal maps are only nonexpansive when marginal entries
vanish.  A fully formal statement must specify exactly which four vectors are strictly positive and
which may inherit zeros from $p$ or $q$.
\item The source proof contains a typographical cross-reference to the same discrete-system equation
twice where two different equations are intended.  The mathematical use is the pair of product
constraints for $p$ and $q$.
\item Projective convergence alone does not choose a norm normalization.  The paper's conclusion
that the raw iteration converges to a vector is tied to the homogeneous cycle and the mass argument;
a detailed analytic treatment should make the normalization control explicit.
\end{commentarynotes}

\begin{thebibliography}{99}
\bibitem{CGP2021}
Y. Chen, T. T. Georgiou, and M. Pavon,
\emph{Stochastic Control Liaisons: Richard Sinkhorn Meets Gaspard Monge on a Schr\"odinger Bridge},
SIAM Review 63(2):249--313, 2021.

\bibitem{Girsanov1960}
I. V. Girsanov,
``On transforming a certain class of stochastic processes by absolutely continuous substitution of measures,''
\emph{Theory of Probability and Its Applications} 5(3):285--301, 1960.

\bibitem{KaratzasShreve1988}
I. Karatzas and S. E. Shreve,
\emph{Brownian Motion and Stochastic Calculus},
Springer, 1988.

\bibitem{Fortet1940}
R. Fortet,
``R\'esolution d'un syst\`eme d'\'equations de M. Schr\"odinger,''
\emph{Journal de Math\'ematiques Pures et Appliqu\'ees}, ninth series, 19:83--105, 1940.

\bibitem{Beurling1960}
A. Beurling,
``An automorphism of product measures,''
\emph{Annals of Mathematics} 72:189--200, 1960.

\bibitem{Jamison1974}
B. Jamison,
``Reciprocal processes,''
\emph{Zeitschrift f\"ur Wahrscheinlichkeitstheorie und Verwandte Gebiete} 30:65--86, 1974.

\bibitem{Jamison1975}
B. Jamison,
``The Markov processes of Schr\"odinger,''
\emph{Zeitschrift f\"ur Wahrscheinlichkeitstheorie und Verwandte Gebiete} 32:323--331, 1975.

\bibitem{Follmer1988}
H. F\"ollmer,
``Random fields and diffusion processes,''
in \emph{\'Ecole d'\'Et\'e de Probabilit\'es de Saint-Flour XV--XVII},
Lecture Notes in Mathematics 1362, Springer, 1988.

\bibitem{Birkhoff1957}
G. Birkhoff,
``Extensions of Jentzsch's theorem,''
\emph{Transactions of the American Mathematical Society} 85(1):219--227, 1957.

\bibitem{Bushell1973Projective}
P. J. Bushell,
``On the projective contraction ratio for positive linear mappings,''
\emph{Journal of the London Mathematical Society} 2:256--258, 1973.

\bibitem{Bushell1973Hilbert}
P. J. Bushell,
``Hilbert's metric and positive contraction mappings in a Banach space,''
\emph{Archive for Rational Mechanics and Analysis} 52:330--338, 1973.

\bibitem{Sinkhorn1964}
R. Sinkhorn,
``A relationship between arbitrary positive matrices and doubly stochastic matrices,''
\emph{Annals of Mathematical Statistics} 35(2):876--879, 1964.

\bibitem{SinkhornKnopp1967}
R. Sinkhorn and P. Knopp,
``Concerning nonnegative matrices and doubly stochastic matrices,''
\emph{Pacific Journal of Mathematics} 21(2):343--348, 1967.

\bibitem{FranklinLorenz1989}
J. Franklin and J. Lorenz,
``On the scaling of multidimensional matrices,''
\emph{Linear Algebra and its Applications} 114/115:717--735, 1989.

\bibitem{EssidPavon2019}
M. Essid and M. Pavon,
``Traversing the Schr\"odinger Bridge strait: Robert Fortet's marvelous proof redux,''
\emph{Journal of Optimization Theory and Applications} 181(1):23--60, 2019.

\bibitem{Leonard2019}
C. L\'eonard,
``Revisiting Fortet's proof of existence of a solution to the Schr\"odinger system,''
arXiv:1904.13211, 2019.
\end{thebibliography}

\end{document}
